% ===== §8 The Inventory of Strategies and the Teleological Filter =====
\section{The Inventory of Strategies, and the Teleological Filter}
\label{p17:sec:practices}

\noindent
The previous section licensed a political economy of the fence and forbade its migration to the bloom. This section conducts that political economy. It surveys what the long history of opaque actors surviving under anarchy has actually devised---the inventory of IR strategies---and then subjects each to a filter the international case never required: the \emph{teleological filter} that asks, of every strategy, not merely whether it secures the actor but whether the securing damages the end the security was for. Many strategies that serve a state perfectly well destroy a dyad, and the filter is what tells them apart. The section is thus the hinge between the diagnostic chapters and the practical ones: it builds the toolbox the praxis (\xref{p17:sec:praxis}) will use, and marks, on each tool, whether it may be used at all.

\subsection{The classical inventory}
\label{p17:sub:practices-inventory}

The strategies opaque actors have devised for surviving a field of power fall into three families, to which a fourth, institutional set may be added.

The first family is \emph{balancing}---the generation of counter-weight against accumulating power. It subdivides, in the classical literature, into \emph{internal balancing} (the building-up of one's own capacities rather than reliance on others) and \emph{external balancing} (alliance with others against a common power), and into the finer manoeuvres of \emph{buck-passing} (letting another bear the cost of balancing while one free-rides), \emph{chain-ganging} (being bound, through alliance, into another's conflicts), and \emph{offshore balancing} (holding back, intervening only as needed to maintain a balance one does not directly join).\footnote{For the internal/external distinction and the dynamics of balancing under anarchy, see \textcite{waltz1979}; for buck-passing and chain-ganging as recurrent alliance pathologies, the classical literature on alliance behaviour.}

The second family is \emph{accommodation}---survival through yielding rather than opposing. It includes \emph{bandwagoning} (aligning with the stronger party rather than against it, typically through asymmetric concession and the acceptance of a subordinate role), \emph{appeasement} (yielding a part to forestall the loss of the whole), and \emph{compensation} (granting an interest elsewhere to secure one's core concern).

The third family is \emph{avoidance}---survival by declining to enter the contest at all: \emph{neutrality} (refusing to take sides) and \emph{isolation} or \emph{exit} (reducing one's surface of contact with the power-field).

To these the liberal-institutionalist tradition adds a fourth set, aimed not at relative position but at the stabilisation of expectations and the reduction of misperception: the \emph{concert} (management of conflict through regular consultation), \emph{confidence-building measures} (incremental, symmetrical, verifiable steps that build trust between parties who cannot inspect one another's intentions), and \emph{calculated ambiguity} or the maintenance of a buffer (the preservation of strategic latitude so as to avoid being forced to a showdown).

\subsection{The teleological filter}
\label{p17:sub:practices-filter}

Now the filter. A state may use any of these as its situation dictates, for a state's only end is survival and any strategy that secures survival is, by that fact, justified. A dyad cannot, because its survival is a means to a further end---the generativity of the love---and a strategy that secures the means while corrupting the end is, for a dyad, a defeat dressed as a victory. The filter sorts the inventory into three classes.

\begin{claim}[The teleological filter]
A strategy admissible for a state, whose end is survival, may be inadmissible for a dyad, whose survival is merely the means to the generativity it protects. Every strategy must therefore be filtered by whether the securing damages the secured. Three classes result: strategies \emph{directly admissible} (they protect the conditions of autonomy and are teleologically neutral); strategies admissible \emph{only after reconstruction} (sound in form but carrying a relative-gains currency that must be stripped before use); and strategies \emph{inadmissible} (they secure survival precisely by corrupting the end, so that their success is the defeat).
\end{claim}

\paragraph{Directly admissible.} Three of the newer strategies pass the filter unaltered, because their logic is the protection of autonomy rather than the pursuit of relative advantage. \emph{Resilience}---the reduction of exposure to single points of failure, the building of redundancy and diversified channels---is teleologically neutral: it secures a party's autonomy without reckoning in anyone's relative loss, and recent work on small-state survival has rightly shifted the emphasis from balancing toward exactly this.\footnote{On the shift from balancing to resilience in small-state strategy---autonomy as the management of one's own vulnerabilities rather than dependence on external guarantee, and resilience as reduced exposure to single points of failure rather than self-sufficiency---see the contemporary small-state literature.} \emph{Hedging}---the refusal to be forced into a binary alignment, the maintenance of latitude through deliberately ambiguous signalling---protects the bond's space against capture by an external script, and is, as \xref{p17:sub:practices-why} will argue, peculiarly suited to the dyad. And the \emph{buffer} or calculated ambiguity protects a core from premature exposure. These three tend the fence and reckon no relative gain; they are admitted whole.

\paragraph{Admissible only after reconstruction.} \emph{Balancing} and \emph{external alliance} are sound in form---the generation of counter-weight against a power is exactly what the suppression of freezing will require (\xref{p17:sec:balance})---but they carry, in their classical statement, the currency of relative gain, and that currency must be stripped before they cross into intimate use. A dyad may form an alliance against an external power (two partners presenting a united front to a coercive family); but the moment the calculative logic of the alliance is turned \emph{inward}---the moment the partners begin to manage one another as an alliance manages its members, with united fronts and coordinated positions inside the dyad---it has become \emph{regime bleeding} (\xref{p17:sub:ir-gradient}), the periphery's logic installed at the core. Likewise \emph{fait accompli}: legitimate as a move against an external veto-holder, but, turned between the partners themselves, simple manipulation. These strategies are admitted only on condition that their relative-gains currency is stripped and their operation confined to the external face.

\paragraph{Inadmissible.} \emph{Bandwagoning}---survival through asymmetric concession and the acceptance of subordination---is the paradigm of the strategy that secures the means by destroying the end. A partner may indeed secure the relationship's continuation by yielding, accepting the subordinate role, ceasing to dissent; the bond survives. But what survives is no longer a generative coupling, for the love required the mutual subjecthood that the bandwagoning has surrendered. The relationship persists as a structure of domination---which is to say, the very pathology the whole enterprise was meant to prevent. So too any strategy that treats the partner as a relative-gains adversary: it may win the local contest and has, in winning, lost the war. These are not admitted at all.

\subsection{Why the dyad is harder than the state}
\label{p17:sub:practices-why}

The filter exposes, in concrete strategic terms, the greater complexity the paper has claimed for the small system. A state may treat its strategies as separable: it balances against rivals and conducts its internal affairs by other logics entirely, and the two need not interfere. The dyad cannot separate them, for the external power it must balance against---the powerful elder, the clan---is frequently \emph{the very same person} who is also, at the core, a member of the bond's own generative field, someone the partners love (cause 16, the inseparability of identity). One cannot balance cleanly against a beloved. The classical binary choice---\emph{balance or bandwagon}---therefore fails for the dyad, because to balance against the elder is to wound the partner's own bond to their family, while to bandwagon is to surrender the dyad's autonomy. Neither pure strategy is available.

This is the deep structural reason that hedging suits the dyad as no classical strategy does: hedging is precisely the refusal of the binary, the maintenance of a position that is neither balancing nor bandwagoning but a third thing, holding latitude open where the classical choice would force a showdown that the inseparability of identity makes catastrophic. What in the state system is a sophisticated option is, in the dyad, very often the only non-destructive move.

\subsection{The newer resources: power's many dimensions, and the right to justification}
\label{p17:sub:practices-newer}

Two strands of recent scholarship deserve separate notice, because they supply the dyad with resources the classical inventory lacks, and because both bear directly on the suppression of freezing to come.

The first is the recognition that power is \emph{multidimensional}, so that an actor weak along one axis may be strong along another and recover agency through it.\footnote{For the typology of power across compulsory, institutional, structural, and productive dimensions---and thus for the possibility that an actor weak in material terms may be strong in others---see \textcite{barnett2005}; for the everyday, infrapolitical forms of resistance available to the materially subordinate, \textcite{scott1985}.} This matters because the classical inventory, fixated on a single axis of relative capability, cannot see the counter-force that a partner weak in resources may exercise through narrative, through the very dependency of the stronger party's affection, through the infrapolitical ``weapons of the weak.'' The multidimensionality of power is, as \xref{p17:sec:balance} will argue, what makes generative balance possible at all: it is because power has many sources that no single accumulation can be final, and counter-force can always, in principle, be generated along an axis the dominant power does not control.

The second is the \emph{right to justification}: the principle that relational domination consists in one party's being subject to another through an asymmetry that the subordinated party cannot demand be justified to them, and that the corrective is the standing entitlement to require justification.\footnote{For relational domination as subjection to an asymmetry of power that need not be justified to the subordinated, and for the right to justification as its normative corrective, see \textcite{forst2012}.} This furnishes the symbolic order with a concrete counter-power mechanism, internal as well as external: the entitlement of either party to demand that the other's exercise of power be justified to them drags a power that would otherwise operate silently---the hidden power lodged in defaults and unspoken rules---back into the space of the answerable. It is, as the praxis will show (\xref{p17:sub:praxis-inner}), the principle underwriting the practice of meta-communication and of dissent.

\subsection{The empirical caution: balance is not a law of nature}
\label{p17:sub:practices-empirical}

A final empirical finding frames everything the next section will build, and it must be entered here, where the strategies are inventoried, lest the praxis rest on a false comfort. It is tempting to suppose that power, left alone, balances itself---that counter-force arises spontaneously, so that a relation need only avoid interference and equilibrium will return. The historical record refuses this comfort. The systematic study of balancing across many international systems finds no tendency of systems to equilibrate on their own; where hegemony was prevented, the preventing force often lay outside the system rather than arising within it.\footnote{For the finding that international systems do not reliably self-equilibrate, and that the strong checks on hegemony have often lain outside the system rather than within it, see the comparative-historical literature testing balance-of-power theory across world history.} The lesson transfers directly and soberingly to the dyad: counter-force does \emph{not} arise on its own. Power, unopposed, accumulates and freezes; the balance that keeps a relation generative is not a natural equilibrium to be awaited but a force to be continually \emph{generated}. This finding is the empirical ground of the negative summit to which the paper now turns---the account of generative balance, and of why the freezing of power must be actively, perpetually opposed.
